\section{第 13.1 节：氢原子的本征问题}
\label{section:QM-C13-S01}
\studymeta
  {QM-C13-S01}
  {2026-10-24 至 2026-10-25}
  {Shankar, \textit{Principles of Quantum Mechanics}, Second Edition}
  {第 13 章 \S 13.1 The Eigenvalue Problem；印刷 pp.~353--359，至 \S 13.2 前 / PDF pp.~361--367}
  {四块讲解与习题 13.1.3、13.1.5 已核对；笔记已确认}

\textbf{主线。}二体问题分离出内部运动，库仑径向方程的两端条件筛选离散能量；截断级数给出波函数，再由径向概率和维里定理描述原子大小及能量分配。

\textbf{记号与近似。}$m_e$ 是电子质量，$M$ 是质子质量，$m$ 只表示磁量子数；$e>0$ 是电荷量的大小。沿用教材 Gaussian 单位，$V=-e^2/r$；在 SI 中把库仑耦合 $e^2$ 替换为 $e^2/(4\pi\varepsilon_0)$。本节研究无外场的非相对论纯库仑模型，不计自旋及相对论修正；完成二体分离后采用固定质子近似。

\notetopic{QM-C13-S01-T01}{二体分离、内部能量与径向边界}
\keyidea{原子的整体平动与电子、质子的相对运动是不同自由度；氢原子能级在这里指内部 Hamiltonian 的本征值。}
\textbf{教材位置：}印刷 pp.~353--354 上部 / PDF pp.~361--362；质心分离的显式写法连接 \S 2.3 与 \S 10.1。

定义质心与相对坐标
\begin{equation}
  \mathbf R_{\mathrm{cm}}=\frac{m_e\mathbf r_e+M\mathbf r_p}{m_e+M},
  \qquad \mathbf r=\mathbf r_e-\mathbf r_p,\qquad
  \mu=\frac{m_eM}{m_e+M}.
\end{equation}
相应动量为 $\mathbf P_{\mathrm{cm}}=\mathbf p_e+\mathbf p_p$、
$\mathbf p=(M\mathbf p_e-m_e\mathbf p_p)/(m_e+M)$，因而
\begin{equation}
  \boxed{H_{\mathrm{tot}}=
    \frac{\mathbf P_{\mathrm{cm}}^2}{2(m_e+M)}
    +\underbrace{\left(\frac{\mathbf p^2}{2\mu}-\frac{e^2}{r}\right)}_{H_{\mathrm{int}}}.}
\end{equation}
这一分离在上述模型内精确成立；增加整体平动动能不会改变内部束缚能级。教材随后令 $M\to\infty$、$\mu\to m_e$，将质子视为固定；这一步才是近似。

\dialoguesupplement{$\mathbf r$ 是相对位置矢量，$r=|\mathbf r|$ 是电子与质子的距离。固定质子并把它放在原点时，$r$ 才同时是电子到该坐标原点的距离；任意实验室原点不能默认与之等同。}

以下 $H$ 与 $E$ 均指内部问题。沿用中心势分离
\begin{equation}
  \psi_{Elm}(r,\theta,\varphi)=R_{El}(r)Y_l^m(\theta,\varphi)
    =\frac{u_{El}(r)}rY_l^m(\theta,\varphi),
\end{equation}
得到
\begin{equation}
  \boxed{-\frac{\hbar^2}{2m_e}u''
    +\left[\frac{\hbar^2l(l+1)}{2m_er^2}-\frac{e^2}{r}\right]u=Eu.}
  \label{eq:QM-C13-S01-radial}
\end{equation}
离心项仍依赖 $l$，径向方程不含 $m$；此时只能由旋转对称性断言同一径向态、同一 $l$ 下的 $m$ 简并。

取无穷远势能为零。束缚态 $E<0$，结合能 $W=-E>0$，令
\begin{equation}
  \kappa=\frac{\sqrt{2m_eW}}\hbar.
\end{equation}
通常的原点正则条件给出 $u\sim r^{l+1}$；无穷远保留以 $\ee^{-\kappa r}$ 为主要指数因子的衰减支，允许额外幂因子。两端条件及 $\int_0^\infty|u|^2\dd r=1$ 必须同时满足，不能只看到一个衰减指数就宣布归一化成立。

\notetopic{QM-C13-S01-T02}{幂级数截断、主量子数与能级简并}
\keyidea{能量量子化来自原点正则解在无穷远也必须可归一化；不是先规定整数能量，再寻找符合它的波函数。}
\textbf{教材位置：}印刷 pp.~354--356 上部 / PDF pp.~362--364。

定义无量纲变量与参数
\begin{equation}
  \rho=\kappa r,\qquad
  \gamma=\frac{2m_e e^2}{\hbar^2\kappa}
    =e^2\sqrt{\frac{2m_e}{\hbar^2W}}.
\end{equation}
$\gamma$ 对应教材的 $e^2\lambda$。式~\eqref{eq:QM-C13-S01-radial} 变成
\begin{equation}
  \frac{\dd^2u}{\dd\rho^2}
   +\left[-1+\frac\gamma\rho-\frac{l(l+1)}{\rho^2}\right]u=0.
\end{equation}
提出两端的已知因子，令 $u=\rho^{l+1}\ee^{-\rho}F(\rho)$。代入并约去公共因子：
\begin{equation}
  \rho F''+[2(l+1)-2\rho]F'+[\gamma-2(l+1)]F=0,
\end{equation}
此处撇号对 $\rho$ 求导。设 $F=\sum_{k=0}^\infty C_k\rho^k$，比较 $\rho^k$ 系数：
\begin{align}
  &(k+1)(k+2l+2)C_{k+1}
     +[\gamma-2(k+l+1)]C_k=0,\notag\\
  &\boxed{C_{k+1}=\frac{2(k+l+1)-\gamma}{(k+1)(k+2l+2)}C_k.}
  \label{eq:QM-C13-S01-recurrence}
\end{align}
非零正则解由 $C_0\ne0$ 决定。若级数不截断，大阶系数之比趋于 $2/k$；教材据此作渐近分析，正则解会含 $\ee^{2\rho}$ 型增长，乘上外面的 $\ee^{-\rho}$ 后仍不可归一化。这是大阶级数的渐近判断，不是将有限阶系数逐项等同于指数级数。

令最高非零项为 $C_{n_r}\rho^{n_r}$。要使 $C_{n_r+1}=0$，必须有
\begin{equation}
  \gamma=2(n_r+l+1),\qquad n_r=0,1,2,\ldots.
\end{equation}
定义 principal quantum number（主量子数）$n=n_r+l+1$，则
\begin{equation}
  \kappa=\frac{m_e e^2}{\hbar^2n},\qquad
  \boxed{E_n=-\frac{m_e e^4}{2\hbar^2n^2}
    =-\frac{\mathrm{Ry}}{n^2},\quad
    \mathrm{Ry}=\frac{m_e e^4}{2\hbar^2},\quad n=1,2,\ldots.}
\end{equation}
其中 $\mathrm{Ry}$ 称为里德伯能量（Rydberg energy）。$E_n$ 是负的束缚能量，$W_n=-E_n$ 才是正的电离所需能量；$n$ 增大时能级趋近连续谱阈值 $E=0$，并非等间距。

由 $n_r=n-l-1\geq0$，量子数范围为
\begin{equation}
  \boxed{l=0,1,\ldots,n-1,\qquad m=-l,-l+1,\ldots,l.}
\end{equation}
每个 $l$ 有 $2l+1$ 个 $m$ 态，所以不计自旋时
\begin{equation}
  \boxed{g_n=\sum_{l=0}^{n-1}(2l+1)=n^2.}
\end{equation}
光谱记号 $s,p,d,f$ 分别表示 $l=0,1,2,3$，前面的数字表示 $n$，并未指定 $m$：
\begin{center}
\begin{tabular}{@{}cccc@{}}
  \toprule
  $n$ & 允许的 $l$ & 光谱记号 & 简并度（不计自旋） \\
  \midrule
  $1$ & $0$ & $1s$ & $1$ \\
  $2$ & $0,1$ & $2s,2p$ & $1+3=4$ \\
  $3$ & $0,1,2$ & $3s,3p,3d$ & $1+3+5=9$ \\
  \bottomrule
\end{tabular}
\end{center}
同一 $l$ 下的 $m$ 简并来自旋转对称性；同一 $n$ 下不同 $l$ 的简并是纯库仑问题的额外性质，原因留到 \S 13.2，不能把它推广到任意中心势。

\dialoguesupplement{$E$、$L^2$、$L_z$ 分别由 $n$、$l$、$m$ 标记。固定 $n$ 的任意叠加仍具有确定能量，但若混合不同 $l$，一般没有确定 $L^2$；只混合同一 $l$ 下的不同 $m$，$L^2$ 仍确定为 $l(l+1)\hbar^2$。因此不能以“$m$ 不确定”作为“$L^2$ 不确定”的理由。}

\notetopic{QM-C13-S01-T03}{玻尔半径与氢原子波函数}
\keyidea{球谐函数解决角向部分；径向函数由原点幂因子、衰减指数和截断多项式组成。特殊函数的名称只是对已经得到的多项式作统一记号。}
\textbf{教材位置：}印刷 pp.~356 下部--357 / PDF pp.~364--365。

定义 Bohr radius（玻尔半径）
\begin{equation}
  \boxed{a_0=\frac{\hbar^2}{m_e e^2},\qquad \rho=\frac{r}{na_0}.}
\end{equation}
$a_0$ 是自然长度尺度，不表示电子沿确定的圆形轨道运动。由于 $R=u/r$，原点幂次由 $l+1$ 降为 $l$：
\begin{equation}
  R_{nl}\propto\rho^l\ee^{-\rho}F_{nl}(\rho),
  \qquad \deg F_{nl}=n-l-1.
\end{equation}
这个截断多项式可记为 associated Laguerre polynomial（关联拉盖尔多项式），从而
\begin{equation}
  \boxed{\psi_{nlm}=R_{nl}(r)Y_l^m(\theta,\varphi),\qquad
  R_{nl}(r)\propto\ee^{-r/(na_0)}
    \left(\frac{r}{na_0}\right)^l
    L_{n-l-1}^{\,2l+1}\!\left(\frac{2r}{na_0}\right).}
\end{equation}
此式使用正比号，整体归一化因子尚未指定，不受多项式整体常数约定的影响。角向已归一化时，径向条件是
\begin{equation}
  \int_0^\infty|R_{nl}(r)|^2r^2\dd r=1.
\end{equation}
不能漏掉 $r^2$，也不能把归一化后的 $R$ 误当作完整三维波函数。

最低几个态展示了同能与同形状的区别：
\begin{equation}
  \psi_{100}=\frac{\ee^{-r/a_0}}{\sqrt{\pi a_0^3}},\qquad
  R_{20}\propto\left(2-\frac r{a_0}\right)\ee^{-r/(2a_0)},\qquad
  R_{21}\propto r\ee^{-r/(2a_0)}.
\end{equation}
$2s$ 在 $r=2a_0$ 有一个径向节点，$2p$ 在 $r>0$ 没有径向节点；原点处因 $r^l$ 而取零不计入这些正半轴节点。完整的 $\psi_{210}$ 递推与归一化见\relatedexercise{QM-C13-S01-E01}{教材习题 13.1.3}。

多项式最高幂为 $r^{n-l-1}$，故固定 $n$ 的远处行为为
\begin{equation}
  \boxed{R_{nl}(r)\sim\text{常数}\times r^{n-1}\ee^{-r/(na_0)}
    \qquad(r\to\infty).}
\end{equation}
它的领先形式与 $l$ 无关，但不是说所有 $l$ 的完整径向函数相同。只有 $l=n-1$、多项式为常数时，这一形式才对所有 $r$ 成立。指数中的 $na_0$ 也不能直接读成概率分布的峰值位置。

\notetopic{QM-C13-S01-T04}{径向概率、原子大小与维里定理}
\keyidea{同一个距离 $r$ 可以对应不同测度下的密度；最可能半径、平均半径与平均倒数必须分别计算。}
\textbf{教材位置：}印刷 pp.~358--359 / PDF pp.~366--367；基态两种密度的对照是学习对话补充。

\textbf{1. 从单位体积到球壳。}一般的径向概率密度为
\begin{equation}
  \boxed{p_r(r)=r^2\int|\psi(r,\theta,\varphi)|^2\dd\Omega,
    \qquad \Pr(r\in[a,b])=\int_a^b p_r(r)\dd r.}
\end{equation}
对角向已归一化的乘积态，$p_r=r^2|R_{nl}|^2=|u_{nl}|^2$。只有完整波函数球对称时，才能直接写成 $p_r=4\pi r^2|\psi(r)|^2$。位置密度 $|\psi|^2$ 的量纲是长度的 $-3$ 次方，$p_r$ 的量纲是长度的 $-1$ 次方。

基态中
\begin{equation}
  |\psi_{100}|^2=\frac{\ee^{-2r/a_0}}{\pi a_0^3},\qquad
  p_r(r)=\frac{4r^2}{a_0^3}\ee^{-2r/a_0}.
\end{equation}
前者在原点最大，后者在原点为零、在 $r=a_0$ 最大；相同厚度球壳的体积因子 $4\pi r^2$ 改变了分布。这不是把相对距离换成了另一种距离，也不表示连续变量在单个精确半径上有非零概率。

\begin{center}
\begin{minipage}{\linewidth}
\centering
\begin{tikzpicture}[x=1.15cm,y=1.8cm,>=Stealth]
  \begin{scope}
    \draw[->] (0,0)--(4.4,0) node[right] {$r/a_0$};
    \draw[->] (0,0)--(0,1.25);
    \node[above] at (1.8,1.25) {$\pi a_0^3|\psi_{100}|^2$};
    \draw[thick,noteBlue,domain=0:4,samples=90] plot (\x,{exp(-2*\x)});
    \node[left] at (0,1) {$1$};
    \node[below left] at (0,0) {$0$};
    \node[below] at (1,0) {$1$};
  \end{scope}
  \begin{scope}[shift={(6.2,0)}]
    \draw[->] (0,0)--(4.4,0) node[right] {$r/a_0$};
    \draw[->] (0,0)--(0,1.25);
    \node[above] at (1.8,1.25) {$a_0p_r(r)$};
    \draw[thick,accentRed,domain=0:4,samples=90] plot (\x,{4*\x*\x*exp(-2*\x)});
    \draw[dashed,gray] (1,0)--(1,{4*exp(-2)});
    \node[above] at (1,{4*exp(-2)}) {峰值};
    \node[below left] at (0,0) {$0$};
    \node[below] at (1,0) {$1$};
  \end{scope}
\end{tikzpicture}

\small 按上述解析式原创绘制；两图纵轴分别作无量纲化，不是同一个概率密度。
\end{minipage}
\end{center}

\textbf{2. 最可能半径与平均半径。}先取 $l=n-1$，此时 $R\propto r^{n-1}\ee^{-r/(na_0)}$，故
\begin{equation}
  p_r\propto r^{2n}\ee^{-2r/(na_0)},\qquad
  \frac{\dd\ln p_r}{\dd r}=\frac{2n}r-\frac2{na_0}=0
  \quad\Longrightarrow\quad
  \boxed{r_{\mathrm{peak}}=n^2a_0\quad(l=n-1).}
\end{equation}
不能把这个峰值公式不加条件地用于其他 $l$。一般态有 $n-l-1$ 个正半轴径向节点，径向概率密度可有多个峰。教材直接给出、未展开积分证明的平均半径公式为
\begin{equation}
  \boxed{\expect r_{nlm}=\frac{a_0}{2}[3n^2-l(l+1)].}
\end{equation}
它说明平均半径依赖 $n,l$，不依赖 $m$，大小尺度为 $n^2a_0$。例如基态 $r_{\mathrm{peak}}=a_0$，但 $\expect r=3a_0/2$；$2s$ 与 $2p$ 虽同能，平均半径分别为 $6a_0$ 与 $5a_0$。

\textbf{3. 能量与平均倒数。}对库仑束缚能量本征态，virial theorem（维里定理）给出
\begin{equation}
  \boxed{2\expect T=-\expect V,\qquad
    \expect T=-E_n,\qquad \expect V=2E_n.}
\end{equation}
其对易子证明见\relatedexercise{QM-C13-S01-E02}{教材习题 13.1.5}。这是定态的期望值关系，不是 $2T+V=0$ 的算符恒等式，也不是任意波包任意时刻都成立的条件。

将 $V=-e^2/r$ 与能谱合用：
\begin{equation}
  E_n=\frac12\expect V=-\frac{e^2}{2}\expect{\frac1r}
    =-\frac{e^2}{2a_0n^2}
  \quad\Longrightarrow\quad
  \boxed{\expect{\frac1r}=\frac1{a_0n^2}.}
\end{equation}
同一 $n$ 的所有 $l,m$ 态有相同的平均倒数，却不必有相同平均半径。尤其
\begin{equation}
  \expect{\frac1r}\ne\frac1{\expect r}\quad\text{一般成立；基态有}\quad
  \expect{\frac1r}=\frac1{a_0},\qquad \frac1{\expect r}=\frac2{3a_0}.
\end{equation}
教材由平均倒数提示原子的大小尺度；不能仅把它取倒数就宣称得到精确 $\expect r$。精确平均半径应使用前面的公式，非线性函数与取期望不能普遍交换。

\textbf{一分钟回顾。}质心平动不改变本模型的内部能级；两端条件使径向级数截断，得到 $n=n_r+l+1$、$E_n=-\mathrm{Ry}/n^2$ 及不计自旋的 $n^2$ 重简并。波函数是径向因子乘球谐函数，$l$ 决定 $L^2$、$m$ 决定 $L_z$。原子大小要从 $p_r=r^2|R|^2$ 判断，不能混淆体积密度、峰值与平均值；维里定理联系定态的动能、势能和 $\expect{1/r}$。

\section{第 13.2 节：氢原子能谱的简并}
\label{section:QM-C13-S02}
\studymeta
  {QM-C13-S02}
  {2026-10-25}
  {Shankar, \textit{Principles of Quantum Mechanics}, Second Edition}
  {第 13 章 \S 13.2 The Degeneracy of the Hydrogen Spectrum；印刷 pp.~359 底--361 上部，至 \S 13.3 前 / PDF pp.~367--369}
  {两块讲解与教材习题 13.2.1 已核对；笔记已确认}

\textbf{主线。}旋转对称性解释同一 $l$ 内的 $m$ 简并；库仑问题不同 $l$ 之间的额外简并与 Runge--Lenz 守恒量有关。本节指出这一联系，完整的量子对称性论证留到 \S 15.4。

\textbf{模型与记号。}沿用固定质子、无外场、不计自旋及相对论修正的纯库仑模型和 Gaussian 单位。$m_e$ 表示电子质量，$m$ 表示磁量子数；经典矢量 $\mathbf l,\mathbf n$ 与量子算符 $\mathbf L,\mathbf N$ 分开书写。粗体 $\mathbf n$ 不是主量子数 $n$。

\notetopic{QM-C13-S02-T01}{旋转对称性与两层简并}
\keyidea{能量不依赖 $m$ 与不依赖 $l$ 是两件事；普通旋转可以改变方向，却不改变角动量量子数 $l$。}
\textbf{教材位置：}印刷 pp.~359 底--360 / PDF pp.~367--368。

氢原子能量 $E_n=-\mathrm{Ry}/n^2$ 只依赖 $n$。例如 $n=2$ 的四个轨道态为
\begin{equation}
  \underbrace{\ket{2,0,0}}_{2s},\qquad
  \underbrace{\ket{2,1,-1},\ \ket{2,1,0},\ \ket{2,1,1}}_{2p}.
\end{equation}
三个 $2p$ 态彼此等能，与 $2s$ 也同能，属于两层不同的简并。

对中心势，旋转不变性给出 $[H,L_i]=0$，从而 $[H,L_\pm]=0$。若 $H\ket{n,l,m}=E_n\ket{n,l,m}$，则
\begin{equation}
  H\bigl(L_\pm\ket{n,l,m}\bigr)
    =L_\pm H\ket{n,l,m}
    =E_n\bigl(L_\pm\ket{n,l,m}\bigr).
\end{equation}
只要升降后的向量非零，它仍具有同一能量；这解释固定 $l$ 内的 $m$ 简并。但
\begin{equation}
  [L^2,L_\pm]=0,\qquad
  L^2\bigl(L_\pm\ket{n,l,m}\bigr)
    =l(l+1)\hbar^2\bigl(L_\pm\ket{n,l,m}\bigr),
\end{equation}
说明升降不改变 $l$。一般的有限旋转同样与 $L^2$ 对易，不能将 $2s$ 旋转成 $2p$。

对任意中心势 $V(r)$，旋转对称性只保证固定径向态、固定 $l$ 内的 $m$ 简并，并不强制不同 $l$ 等能。教材称库仑问题的后一种简并为 accidental degeneracy（偶然简并），引号强调它并非没有原因，而是普通旋转对称性尚未解释的额外简并。

\notetopic{QM-C13-S02-T02}{Runge--Lenz 矢量与额外守恒量}
\keyidea{库仑系统不仅有固定轨道平面的角动量，还具有固定椭圆朝向的额外守恒矢量；其量子对应提供联系不同 $l$ 同能态的结构。}
\textbf{教材位置：}印刷 p.~360，式 (13.2.1)--(13.2.3) 及脚注；几何解释与习题延续至 p.~361 上部 / PDF pp.~368--369。

\textbf{1. 经典轨道中固定什么？}在 $V=-e^2/r$ 中，教材给出 Runge--Lenz vector（龙格--楞次矢量）
\begin{equation}
  \boxed{\mathbf n=\frac{\mathbf p\times\mathbf l}{m_e}
     -e^2\frac{\mathbf r}{r},\qquad
     \mathbf l=\mathbf r\times\mathbf p,\qquad
     \frac{\dd\mathbf n}{\dd t}=0.}
\end{equation}
对角动量非零的非圆束缚轨道，$\mathbf l$ 固定轨道平面，$\mathbf n$ 位于该平面内，沿椭圆长轴指向最近点。$\mathbf n$ 的方向不变，远、近点的朝向便不会逐圈转动，即无近心点进动。一般中心势的角动量守恒只保证平面固定，并不保证轨道闭合。圆轨道则有 $\mathbf n=0$，没有特别的近心点方向。

\begin{center}
\begin{minipage}{\linewidth}
\centering
\begin{tikzpicture}[scale=0.9,>=Stealth]
  \draw[thick,noteBlue] (-1,0) ellipse (2 and 1.73205);
  \draw[dashed,gray] (-3.25,0)--(1.25,0);
  \draw[->,thick] (0,0)--(-3,0) node[midway,below] {$\mathbf r_{\max}$};
  \draw[->,thick] (0,0)--(1,0) node[midway,below] {$\mathbf r_{\min}$};
  \fill (0,0) circle (2pt) node[below=6pt] {力心};
  \fill (-3,0) circle (1.5pt) node[left] {最远点};
  \fill (1,0) circle (1.5pt) node[right] {最近点};
  \draw[->,very thick,accentRed] (-0.1,0.75)--(1,0.75)
    node[midway,above] {$\mathbf n$ 的方向};
\end{tikzpicture}

\small 原创示意：力心位于椭圆的一个焦点；红箭头只示方向，不表示位置矢量。
\end{minipage}
\end{center}

转折点关系、远近方向及圆轨道的推导见\relatedexercise{QM-C13-S02-E01}{教材习题 13.2.1}。这里讨论的是经典轨道，不能据此给量子能量本征态指定一条电子运行轨迹。

\textbf{2. 量子化为何需要对称化？}由于 $\mathbf P$ 与 $\mathbf L$ 的分量一般不对易，不能机械地把经典叉积照搬为算符。按分量取伴随，有
\begin{equation}
  (\mathbf P\times\mathbf L)^\dagger=-\mathbf L\times\mathbf P.
\end{equation}
教材取该算符与其伴随的平均，得到
\begin{equation}
  \boxed{\mathbf N=\frac{\mathbf P\times\mathbf L-\mathbf L\times\mathbf P}{2m_e}
       -e^2\frac{\mathbf R}{R},\qquad R=\sqrt{X^2+Y^2+Z^2}.}
\end{equation}
前面虽写成相减，实际正是对称化，保证分量的形式厄米性。在经典极限中，$\mathbf l\times\mathbf p=-\mathbf p\times\mathbf l$，它恢复为原来的 $\mathbf p\times\mathbf l/m_e$。

\textbf{3. 守恒量如何联系同能态？}对纯库仑 Hamiltonian
\begin{equation}
  H=\frac{\mathbf P^2}{2m_e}-\frac{e^2}{R},\qquad
  \boxed{[N_i,H]=0\quad(i=x,y,z).}
\end{equation}
教材在本节直接给出这一对易关系，未展开逐项计算。若 $H\ket\psi=E\ket\psi$，则
\begin{equation}
  H(N_i\ket\psi)=N_iH\ket\psi=E(N_i\ket\psi).
\end{equation}
所以 $N_i\ket\psi$ 非零时仍属于同一能量本征空间。与保持 $l$ 的普通旋转不同，利用 $\mathbf N$ 可构造联系不同 $l$ 同能态的算符；具体构造及其对称性解释留到 \S 15.4，不把本节的线索当作已完成的完整证明。

\dialoguesupplement{仅有 $[Q,H]=0$ 不能直接证明额外简并：还要证明 $Q\ket\psi$ 非零且与原态线性无关。例如 $Q=H$ 只给能量本征态乘以 $E$，没有产生新态。}

\textbf{一分钟回顾。}固定 $l$ 的 $m$ 简并来自旋转对称性；不同 $l$ 的库仑简并需要额外对称性。经典 $\mathbf n$ 沿椭圆长轴指向近点，圆轨道中为零；量子 $\mathbf N$ 必须处理算符次序，并与库仑 $H$ 对易。同能映射是否产生独立新态是证明简并的必要检查，完整联系在 \S 15.4 展开。
